% Chapter~\ref{ch:anatomy_joint_modeling}: Anatomy and Joint Modeling
% The Physics of Golf: Force, Drift, and Control in the Golf Swing

\chapter[Anatomy and Joint Modeling: Choosing the Right Idealization]{Anatomy and Joint Modeling:\\Choosing the Right Idealization}
\label{ch:anatomy_joint_modeling}

\begin{laymansbox}{What a Joint Model Promises}{box:ch22_modeling_problem}
A knee model can answer several different questions: where the foot moves, what net moment acts between the leg segments, how cartilage carries contact load, or which tissue might be injured. A model adequate for the first question may be inadequate for the others. The useful simplification is therefore defined by the prediction and the evidence available to test it.

Joint primitives give us a precise language for allowed relative motion. They are one part of a biomechanical model, alongside muscles, passive tissues, segment inertia, grip and ground contact. This chapter connects those choices: how a change in anatomy or constraints propagates into motion, force, work and the conclusions we can defend about the golf swing.
\end{laymansbox}

\section{The Modeling Problem: Choose the Question First}\label{sec-modeling_problem}\label{sec:modeling_problem}

Consider the knee. A rigid-body model may replace the tibiofemoral articulation with a hinge. A more detailed model can include coupled rotations and translations, patellar motion, ligament forces and contact between deformable surfaces. These models represent different aspects of the same structure. Neither the number of coordinates nor the realism of a rendered skeleton establishes predictive accuracy.

A hinge has one configuration coordinate, its relative angle \(q\). That angle is not automatically an actuator input. In forward dynamics, a net moment or a set of muscle excitations can drive the motion; the angle evolves from the initial state and the equations of motion. In a prescribed-motion calculation the angle trajectory is supplied, and inverse dynamics can estimate the net forces needed to produce it. Identifying individual muscle forces or tissue stress requires additional assumptions and observations.

For a fixed hinge axis \(u\) expressed in the parent joint frame, one possible orientation convention is

\[
R(q)=\exp(q[u]_\times)R_0,
\]

where \(u\) is a unit vector, \([u]_\times v=u\times v\), and \(R_0\) is the relative orientation at \(q=0\). This formula describes rotation about the specified axis. The joint-origin constraint and the placement of each segment's center of mass are also needed to determine the segment motion.

A detailed finite-element model can resolve deformation and contact effects that this hinge removes, provided its geometry, materials, boundary conditions and numerical approximation are supported. It does not become accurate merely by adding detail. Conversely, a reduced model is not automatically identifiable or fast enough for a particular optimizer. Calibration, sensitivity analysis, convergence checks and validation against independent observations establish which predictions are trustworthy.

For golf, begin by declaring the output and system boundary. Club delivery requires the coupled body--hand--club chain and its support conditions. A knee contact-force question also requires the forces that realize the measured motion. A tissue-injury question adds constitutive behavior, exposure and clinical evidence. The following sections develop useful reductions and show where one question requires information that another can omit.
\section{Joint Types in Robotics --- A Taxonomy}\label{sec-joint_types}\label{sec:joint_types}

A joint primitive specifies admissible relative motion. It does not specify the actuator, muscle recruitment, friction, contact pressure or tissue tolerance. For two otherwise free rigid bodies, relative configuration has six dimensions. A regular joint with \(f\) local freedoms imposes \(6-f\) independent constraints. Count the rank of the constraints, not the number of entries in a rotation matrix \citep{Lynch2017}.

In this section, \(R\in SO(3)\) maps vector components from body B's joint frame into body A's joint frame. The vector \(p\) is the displacement from A's joint-frame origin to B's, expressed in A. Attach those origins at the modeled joint center, not at the segment centers of mass. A statement that \(p=0\) therefore permits a segment center of mass to move on an arc. Joint axes are unit vectors and all angular velocities below are relative to A and expressed in A.

\subsection{Revolute (1 DOF)}\label{revolute-1-dof}

A revolute joint permits rotation about one common axis. Let \(u_A\) and \(u_B\) be its axis directions fixed in the respective bodies. Its constraints are

\[
p=0,\qquad R u_B=u_A.
\]

The first equation gives three translational constraints. The equality of two unit axis directions gives two independent rotational constraints: its three written components have differential rank two at a compatible orientation. The total is five, leaving one freedom. With coincident, aligned reference frames and the axis denoted by \(u\), one convenient parameterization is \(R(\theta)=\exp(\theta[u]_\times)\). A general reference orientation introduces a constant rotation factor.

For an arbitrary segment-fixed point with offset \(a\) from the joint, its relative position is \(Ra\). Thus a hinge does not prevent every point of one body from translating relative to the other. It prevents separation of the chosen joint origins. A knee hinge is a possible reduction of biological motion; its adequacy must be assessed for the target output rather than inferred from the primitive's name.

\subsection{Prismatic (1 DOF)}\label{prismatic-1-dof}

A prismatic joint permits sliding along a unit axis \(u\). Choose orthonormal \(e_1,e_2\) perpendicular to \(u\). In aligned reference frames,

\[
\begin{gathered}R=I_3,\qquad p=d u,\\ e_1^T p=0,\qquad e_2^T p=0.\end{gathered}
\]

Fixed orientation imposes three independent conditions on \(SO(3)\), not nine. The two scalar position conditions remove both transverse translations, leaving the distance \(d\). One unspecified perpendicular vector would remove only one transverse direction. A sliding rail is a literal engineering example; a translation coordinate can also be useful in a biological model without implying that an anatomical joint is a perfect rail.

\subsection{Spherical (3 DOF)}\label{spherical-3-dof}

A spherical joint imposes \(p=0\) and permits any \(R\in SO(3)\) before anatomical limits or other loads are added. It has three rotational freedoms. A unit quaternion uses four components with one normalization constraint, and the two signs represent the same orientation; it does not add a fourth rotational freedom.

One local representation is the ZYX sequence

\[
R=R_z(\phi)R_y(\theta)R_x(\psi).
\]

The rotations act about successive axes, whose spatial directions depend on preceding rotations. They are not three fixed, mutually orthogonal axes throughout the motion. At \(\theta=\pm\pi/2\), this Euler coordinate chart loses rank. The ideal spherical joint still has three physical rotational freedoms. A physical three-axis gimbal, in contrast, has a mechanism Jacobian whose rank can fall when its outer and inner axes coincide. A coordinate singularity, a mechanism singularity, a joint limit and loss of task controllability require different analyses.

The hip is often reduced to a spherical joint with an admissible configuration envelope. The primitive's unrestricted orientation set is not a claim that a biological hip can attain every orientation.

\subsection{Universal (2 DOF)}\label{universal-2-dof}

An ideal universal joint consists of two intersecting revolute axes held perpendicular by the intermediate member. Let \(u_A\) be the first axis fixed in A and \(u_B\) the second axis fixed in B. Its finite constraints are

\[
p=0,\qquad u_A^T R u_B=0.
\]

The scalar orthogonality condition supplies one rotational constraint; together with the three position constraints it leaves two freedoms. It does not require a fixed third axis to remain unchanged by \(R\).

For reference axes \(u_A=e_x\), \(u_B=e_y\), use

\[
\begin{gathered}R=R_x(\alpha)R_y(\beta),\\ \omega=W(\alpha)\dot q,\qquad \dot q=(\dot\alpha,\dot\beta)^T,\end{gathered}
\]

\[
\begin{gathered}W(\alpha)=\begin{bmatrix}1&0\\0&\cos\alpha\\0&\sin\alpha\end{bmatrix},\\ W(\alpha)^T W(\alpha)=I_2.\end{gathered}
\]

Both columns remain orthogonal for every \(\alpha\). The ideal two-axis joint therefore does not lock because its own axes align. A driveshaft assembled using a Cardan joint has additional shaft-axis and transmission geometry; its operating limits are not a proof that this two-coordinate joint loses rank. Likewise, a wrist--club task map can be singular even when the local joint rate map has rank two.

The forbidden instantaneous angular-velocity direction is the moving normal \(n=e_x\times R_x(\alpha)e_y\). For example, at \(\alpha=\pi/2\), the allowed second-axis velocity points along A's \(z\) direction. A rule that permanently forbids rotation around A's \(z\) direction would reject legitimate universal-joint motion. At \(\alpha=\pi/2,\beta=0\), \(R e_z=-e_y\), directly disproving the fixed-normal rule.

This rate map also determines which moment components drive the coordinates. If \(m\) is an applied relative moment expressed in A, virtual power gives

\[
Q=W^T m,\qquad Q^T\dot q=m^T\omega.
\]

The generalized moments are projections onto the current axes; they are not arbitrary fixed Cartesian components. An ideal reaction moment along \(n\) does no relative work because \(n^T\omega=0\), although it can transmit load. This distinction matters when interpreting wrist torques, grip reactions and clubhead acceleration.

\subsection{Cylindrical (2 DOF)}\label{cylindrical-2-dof}

A cylindrical joint permits translation and rotation along the same axis. With aligned reference frames, \(p=d u\) and \(Ru=u\) impose two translational and two rotational constraints. The coordinates are \(d\) and \(\theta\). This is a kinematic idealization, not an assertion that one biological joint realizes an exact cylinder.

\subsection{Planar (3 DOF)}\label{planar-3-dof}

A planar joint permits two translations in a plane and rotation about its normal. For a fixed unit normal \(n\), aligned reference normals give \(n^T p=0\) and \(Rn=n\): one translational and two independent rotational constraints. A foot kept flat on a plane can be modeled this way only within that assumed mode. The primitive alone imposes neither friction nor nonnegative normal force. Sticking, slipping, edge contact and lift-off require additional contact laws and feasibility checks.

\subsection{6-DOF (Free, Floating)}\label{dof-free-floating}

A free body has three translational and three rotational freedoms relative to ground. A floating-base model uses those freedoms to represent overall pose; it does not assert absence of gravity, muscle forces or contact reactions. Quaternion storage may require seven position entries for this six-dimensional pose. Position-array length, velocity dimension and actuator count need not agree.

\subsection{Constraint Counting}\label{constraint-counting}

For a spatial linkage with \(N\) bodies including ground and joints whose constraints are independent, the mobility estimate is

\[
d=6(N-1)-\sum_j(6-f_j).
\]

Count three moving links attached in series to a fixed base. If the base is merged with ground, \(N=4\) and there are three revolute joints:

\[
d=6(4-1)-3(6-1)=3.
\]

If the fixed base is kept as a separate body, \(N=5\) and its weld to ground must be counted as a fourth joint with \(f=0\). Then \(d=24-(6+15)=3\). Omitting that weld gives nine freedoms, which describes the floating chain, not the intended fixed-base arm. The error is missing topology, not the use of a spatial formula on a planar example. The formula can be applied to trees as well as loops when its assumptions hold.

For a fixed-base tree, independent joint coordinates give \(d=\sum_j f_j\); a floating base adds six. In a loop, let \(q\in\mathbb R^{n_q}\) be local tree coordinates and \(c(q)=0\) the closure equations. At a regular configuration,

\[
d=n_q-\operatorname{rank}J_c,\qquad J_c=\frac{\partial c}{\partial q}.
\]

Redundant constraints cannot each remove another freedom. At a singular configuration, the linearized tangent space may be larger than the set of realizable finite motions, so a pointwise rank loss must be checked against the nonlinear closure. For a two-handed golf grip, specify the relative hand--club transforms and evaluate that closure rank; do not simply subtract a number called ``the loop.'' Mobility is also distinct from the number of independently controlled inputs and the dimension of a chosen hand or clubhead task.


\section{The Knee --- A Revolute (Mostly)}\label{sec-knee_joint}\label{sec:knee_joint}

\subsection{Anatomy Summary}\label{anatomy-summary}

The tibiofemoral joint carries load between the femoral condyles and tibial plateau through cartilage and the menisci. The menisci distribute contact load and contribute to stability; they are not simply interchangeable shock absorbers. The patella articulates with the femur and changes the geometry of the knee extensor mechanism. A model that merges it into an effective quadriceps moment arm has made a further reduction beyond replacing the tibiofemoral articulation with a hinge.

The cruciate and collateral ligaments resist combinations of displacement and rotation. In a simple anatomical description, the ACL is a major restraint to anterior tibial translation, the PCL to posterior translation, the MCL to valgus loading and the LCL to varus loading. These are load-dependent restraints, not exact kinematic welds. Valgus and varus refer to opposite frontal-plane angulations; they should not be confused with axial twisting. Muscle forces, articular geometry and the capsule also contribute to stability \citep{Neumann2017, Nordin2012}.

\subsection{Range of Motion and the Hinge Approximation}\label{range-of-motion-and-the-hinge-approximation}

Flexion--extension is an important knee motion, but dominance in angular excursion is not sufficient to establish that other motions have negligible effects on a particular output. A hinge may describe a sagittal-plane teaching model well while being unsuitable for estimating axial knee moments, contact loads or how a planted leg accommodates pelvis rotation. Lead and trail knees need not follow the same trajectory. A video projection, an anatomical joint angle and an external knee moment are different measurements.

A commonly quoted flexion range near 140 degrees is an anatomical reference, not the knee excursion used by every golfer \citep{Neumann2017}. Assessment posture, active versus passive testing, subject anatomy and the angle convention matter. An accuracy claim for a golf model needs a declared dataset and error metric, such as hand-position residuals or differences in estimated joint moments.

\subsection{Where Revolute Breaks Down}\label{where-revolute-breaks-down}

Real tibiofemoral motion includes coupled axial rotation, translation and frontal-plane motion. The screw-home description concerns coupling near terminal extension; tibial rotation relative to the femur and femoral rotation relative to a planted tibia must be distinguished. It does not imply an identical rotation in every loading condition or a mechanically irreversible lock.

Comparing an omitted axial rotation with the total flexion range is not an error calculation: these rotations act about different axes. For a point at perpendicular distance \(\rho\) from an omitted rotation axis, an angle error \(\delta\) produces the displacement

\[
\|\Delta p\|=2\rho\left|\sin\frac{\delta}{2}\right|.
\]

A point on the axis has zero displacement from this rotation even though its attached segment changes orientation. A point 5 cm off the axis moves about 8.72 mm for a 10-degree rotation. Tibial length alone is insufficient to calculate the ankle-position error. For club delivery, the relevant effect also depends on the remaining kinematic chain, support constraints and compensating motion.

\subsection{Choosing the Knee Model}\label{choosing-the-knee-model}

Start with the outputs to be explained. For a simplified planar swing, a hinge is a useful explicit assumption. For spatial kinematics, compare it with a model that permits measured coupled rotations or translations, and evaluate residuals over the whole swing. A one-coordinate knee can even have translation and axis orientation prescribed as functions of flexion; one degree of freedom does not require a geometrically fixed hinge.

For tissue loading, additional contact, ligament and muscle models may be necessary, but adding six free coordinates does not itself create a validated injury model. More coordinates require more constraints, constitutive information or measurements to identify their behavior. Computation time must be benchmarked for the actual solver and task; a hinge is not universally required for optimization.

\section{The Hip --- A Ball-and-Socket (Almost)}\label{sec-hip_joint}\label{sec:hip_joint}

\phantomsection\label{the-hip-a-ball-and-socket-almost}{}

\subsection{Anatomy and Spherical Reduction}\label{anatomy-and-spherical-reduction}

The femoral head articulates with the acetabulum. Cartilage, the labrum, capsule, ligaments and surrounding muscles contribute to load transmission and stability. A spherical reduction uses a fixed relative center and three rotational freedoms. Its purpose is to approximate the motion required for a particular analysis, not to assert exact sphericity, zero biological translation or unrestricted motion.

Avoid assigning a universal center shift from an anatomical sketch. A measured shift may reflect real articular motion, a fitted-center definition, marker movement or model mismatch. Its consequence is also task dependent: a center-location error affects segment motion, moment arms and the conversion of a force into a moment. Those effects should be propagated into the quantities being reported.

\subsection{Range of Motion and Golf Measurements}\label{range-of-motion-and-golf-measurements}

Pelvis turn in the laboratory frame is not anatomical hip internal rotation. Hip motion describes the femur relative to the pelvis. Its components depend on anatomical frames, rotation sequence, reference posture and sign convention. Cheetham and colleagues' pelvis--upper-torso separation near transition does not establish a universal lead-hip internal-rotation requirement \citep{Cheetham2001}.

Standard reference values such as approximately 120 degrees of hip flexion describe particular assessments \citep{Neumann2017}. They are not independent limits that can be combined into a rectangular feasible region without checking combined posture. Knee position, loading and active versus passive testing can change the measured range. Measured swing motion must be expressed in the model's frames and coordinates before comparing it with any such limits.

For example, with laboratory orientations \(R_P\) and \(R_F\) for pelvis and femur, respectively, the relative orientation is

\[
R_{PF}=R_P^T R_F.
\]

Changing both laboratory orientations by the same rigid rotation leaves this relative orientation unchanged. Increasing a pelvis-turn angle therefore does not uniquely prescribe a change in hip rotation: the femur and foot may also move. Coordinate components obtained from different Euler sequences are not directly interchangeable.

\subsection{Limits, Impingement and Model Selection}\label{limits-impingement-and-model-selection}

Bony geometry, capsuloligamentous resistance and muscle state can restrict combined hip configurations. A model may describe them through geometric contact, passive forces or an empirically measured admissible envelope. A hard angle bound and a passive resistance law make different predictions near the boundary and should not be used interchangeably without explanation.

Femoroacetabular impingement syndrome is a clinical condition requiring the combination of symptoms, clinical signs and imaging findings; a morphology or a rotation limit alone is not the diagnosis \citep{Griffin2016Warwick}. It is therefore unjustified to assign every affected golfer a 20--30-degree limit or predict a compulsory spinal compensation from that number. Subject-specific observations may motivate a model constraint, whose effects can then be studied without claiming that it diagnoses the cause of pain.

A three-rotation hip with calibrated geometry and configuration limits is a defensible starting point for many whole-body models. Evaluate whether additional translation or contact detail changes the target prediction beyond the uncertainty already present. Motion-capture angles cannot simply be copied into a model with different frames or conventions.

\section{The Shoulder Complex --- A Moving Girdle}\label{sec-shoulder_complex}\label{sec:shoulder_complex}

\subsection{The Four Articulations}\label{the-four-articulations}

The glenohumeral (GH) joint connects humerus to scapula. The acromioclavicular (AC) joint connects scapula to clavicle, and the sternoclavicular (SC) joint connects clavicle to the axial skeleton. The scapulothoracic (ST) articulation describes the scapula's motion against the thorax rather than an additional ordinary synovial joint in series with all three. Its soft-tissue interface and the clavicular linkage constrain how the scapula can move.

This topology matters for counting. One cannot add separately assigned GH, AC, SC and ST freedoms and then announce an approximate independent total. In a model with thorax--clavicle--scapula--humerus segments, ST pose is related to the SC and AC transformations. Additional scapulothoracic contact or coupling equations further restrict the motion. An alternative reduced model may parameterize scapular motion directly, but its independent coordinates and omitted structures must be stated.

\subsection{Moving Centers and Measured Angles}\label{moving-centers-and-measured-angles}

Let \(r_{GH}(q_s)\) locate the GH center relative to the thorax, where \(q_s\) describes shoulder-girdle configuration, and let \(R_H(q_s,q_g)\) orient the humerus. All components are expressed in the thorax frame; the angular velocity below is relative to that frame. For a humerus-fixed point with offset \(a\),

\[
r=r_{GH}(q_s)+R_H(q_s,q_g)a,
\]

\[
\dot r=\frac{\partial r_{GH}}{\partial q_s}\dot q_s
+\omega_H\times(R_Ha).
\]

The moving-center term remains even if a reduced model fits the humeral orientation exactly. With an attached forearm and hand, their additional relative-motion terms must also be included. This is one way scapular motion changes hand paths and grip geometry: it changes both the location from which the arm moves and its orientation.

Humerothoracic orientation describes humerus relative to thorax; GH orientation describes humerus relative to scapula. Widening a fixed-center model's angle limits may help it fit some observed poses, but does not reconstruct the missing center translation, anatomical moment arms or scapular muscle work. A scapula rotating at a prescribed rate is not a fixed scapula.

\subsection{Scapulohumeral Rhythm}\label{scapulohumeral-rhythm}

Scapulohumeral rhythm describes coordinated scapular and humeral motion under a specified task. It is not a universal constant. Direct bone-pin measurements by McClure and colleagues in eight healthy volunteers found substantial scapular rotations about three axes during arm elevation; their reported mean GH-to-ST ratio was 1.7:1 \citep{McClure2001Scapula}. That study does not supply a phase-by-phase golf-swing law.

The familiar 2:1 ratio can still serve as a declared planar illustration: if total arm elevation is 90 degrees and GH elevation is twice the scapular contribution, the components are 60 and 30 degrees. These angles add only in that simplified common-axis construction. General three-dimensional relative orientations compose as rotation matrices; componentwise angle addition is not valid.

In a golf investigation, measure or estimate the relevant shoulder-girdle motion and propagate its uncertainty. Similar hand positions can be reached through different combinations of trunk, scapular, humeral and elbow motion. That kinematic redundancy does not establish which muscles supplied work or how much load reached the GH joint.

\subsection{Choosing the Shoulder Model}\label{choosing-the-shoulder-model}

A fixed-center spherical surrogate can be adequate for a deliberately coarse whole-arm task, provided its limitations are tested against observed hand positions and orientations. Label its angles as effective humerothoracic coordinates when that is what they represent, rather than reporting them as anatomical GH rotations.

For scapular or muscle-work questions, use an explicitly described shoulder-girdle model. Seth and colleagues provide a dedicated scapulothoracic formulation \citep{Seth2016Scapula}. Their later musculoskeletal study illustrates why prescribing scapular motion as a function of humeral motion can alter work attribution: the coupling can allow muscles crossing the GH joint to do work that the model assigns to scapular motion \citep{Seth2019ShoulderWork}. This is a modeling mechanism to examine, not evidence that a particular golf phase is driven by one muscle group.

Their experimental comparison involved one healthy subject performing shrugging and arm-elevation tasks. Muscle work was estimated from the model, with kinematics and EMG used for comparison; muscle forces and velocities were not measured directly. These limits matter when applying its lessons to a golf swing.

Prescribed kinematic coupling also deserves an energy check. If \(q_s=h(q_g)\) is imposed and \(H=\partial h/\partial q_g\), then \(\dot q_s=H\dot q_g\). Generalized forces \(Q_g,Q_s\) reduce to

\[
\begin{gathered}Q_{\rm reduced}=Q_g+H^T Q_s,\\ Q_g^T\dot q_g+Q_s^T\dot q_s=Q_{\rm reduced}^T\dot q_g.\end{gathered}
\]

This identity preserves power for the imposed coupling; it does not establish that the coupling is physiological. Compare models on the same observed task, loads and initial state. A diagnosis or observed weakness does not by itself specify a reduced scapular range, a passive stiffness value or a unique compensatory movement.

\section{The Wrist --- A Universal Joint (Approximately)}\label{sec-wrist_joint}\label{sec:wrist_joint}

\subsection{Anatomy and Primary Motions}\label{anatomy-and-primary-motions}

The radiocarpal and midcarpal articulations involve multiple carpal bones rather than one rigid intermediate member. The familiar wrist coordinates describe net hand motion relative to the forearm: flexion--extension and radial--ulnar deviation. Their axes and centers are approximations that depend on the measurement and model. A two-coordinate universal reduction can be useful without implying that individual carpal bones execute ideal hinges.

Radial deviation means movement toward the thumb side; ulnar deviation means movement toward the little-finger side. Flexion and extension should be defined anatomically, not as ``downward'' and ``upward'' in the laboratory, because forearm orientation changes during the swing. Specify the hand and forearm frames and the order of coordinate rotations before comparing wrist angles across software or studies.

Pronation--supination is primarily forearm rotation involving both the proximal and distal radioulnar articulations. It should not be assigned exclusively to a single elbow-region hinge or counted again as an independent wrist twist. A reduced model may implement it as one forearm rotation coordinate, with an axis appropriate to that model.

\subsection{The Dart-Thrower's Motion}\label{the-dart-throwers-motion}

The dart-thrower's path runs from radial extension toward ulnar flexion. Crisco and colleagues studied its in vivo carpal kinematics and found positions along this path where scaphoid and lunate motion approached zero \citep{Crisco2005Dart}. That finding concerns the measured carpal motions; it does not prove that this is always the strongest or fastest wrist action, or that every golfer should follow that path.

Two-coordinate models can represent a directional stiffness preference. Let \(\Delta q=q-q_0\) be displacement from a chosen reference posture, with coordinates positive in flexion and ulnar deviation. Consider the illustrative elastic law

\[
U=\tfrac12\Delta q^T K\Delta q,\qquad
\tau_{\rm passive}=-K\Delta q,
\]

\[
K=\begin{bmatrix}2&-1\\-1&2\end{bmatrix}
\ \mathrm{N\,m/rad}.
\]

The eigenvalues are 1 and 3 in the stated angular units. Equal positive changes in both coordinates follow the lower-stiffness direction; opposite changes follow the higher-stiffness direction. This is a constructed example of coupling, not a measured human wrist parameter set. Nonlinear stiffness, damping, activation and changing axes can be added when supported by data. The choice of two kinematic coordinates does not forbid coupling between their forces.

\subsection{Wrist Motion, Grip Motion and Club Orientation}\label{wrist-motion-grip-motion-and-club-orientation}

Club roll in space does not imply slip inside the hands. With a rigid hand--club transform,

\[
R_C=R_H R_{HC},\qquad R_{HC}=\text{constant}.
\]

\(R_H\) maps hand-frame components to the laboratory and \(R_{HC}\) maps club-frame components to the hand frame. As the forearm and hand rotate, the club rotates in the laboratory even though it has no relative grip freedom. Add a hand--club roll coordinate only if the model actually permits grip slip, compliance or a bearing-like connection. A shaft-torsion coordinate is another distinct possibility: it describes deformation between locations along the shaft rather than a freely slipping handle.

For an ideal rigid grip, the club's angular velocity equals the hand's angular velocity. Its component along the current shaft direction can vary because of several upstream joint motions. That shaft-axis component, a clubface angle and a chosen Euler roll rate need not be the same observable. With two hands attached to one club, changes at one wrist must also satisfy the other hand's closure constraints; independent angle edits generally break the loop.

\subsection{Choosing the Wrist Model}\label{choosing-the-wrist-model}

Use the two-coordinate reduction for a declared gross-motion question, then test whether its axis geometry, coupled stiffness and grip assumptions explain the measured trajectory and loads. Descriptions such as ``cocking,'' ``bowing'' and ``release'' require a mapping to those coordinates; a projected arm--shaft angle is not a single anatomical wrist angle.

More detailed carpal models can represent motion hidden by a net wrist coordinate, but assigning two independent rotations to both radiocarpal and midcarpal joints is itself a new model assumption, not an anatomical identity. Calibrate the additional coupling and contact information before interpreting individual-bone or tissue loads.

\section{The Elbow --- A Revolute (Very Good Approximation)}\label{sec-elbow_joint}\label{sec:elbow_joint}

\subsection{Anatomy and Two-Coordinate Reduction}\label{anatomy-and-two-coordinate-reduction}

The elbow complex includes humeroulnar, humeroradial and proximal radioulnar articulations. Humeroulnar motion is often reduced to a flexion--extension hinge. Forearm pronation--supination involves the radius and ulna along the forearm, including their distal articulation. A gross upper-limb model can use one flexion coordinate and one forearm rotation coordinate, but two independently specifiable coordinates do not imply dynamically independent motions.

For example, off-diagonal inertia terms allow a moment at one coordinate to change acceleration at another. Muscles can span more than one coordinate, and their moment arms change with geometry. Separating the coordinates in the model is useful bookkeeping; it does not isolate muscle action or energy transfer.

\subsection{The Carrying Angle}\label{the-carrying-angle}

The carrying angle describes the relative alignment of the upper arm and forearm in a specified extended posture, conventionally with the forearm supinated. It is not a complete measurement of the flexion axis. Segment landmark error, forearm rotation and soft-tissue motion affect an estimate. A static angle alone cannot identify a unique three-dimensional axis or muscle moment-arm function.

Estimate segment frames and, where needed, a functional axis from several poses. State which parameters came from anatomical landmarks, fitted motion or imaging. Reference ranges in anatomy texts provide context, but a subject's active range, passive range and task trajectory are different data. Neither a 150-degree flexion reference nor an 80-degree forearm-rotation reference should be treated as a universal swing requirement \citep{Neumann2017}.

\subsection{Choosing the Elbow Model}\label{choosing-the-elbow-model}

A hinge plus a forearm rotation coordinate is a practical starting reduction. Test axis alignment and motion residuals before using it for detailed joint moments. Muscle-force or tendon-load questions require a musculotendon model and additional assumptions; the net elbow moment does not uniquely identify the flexors, extensors or co-contraction.

\section{The Ankle/Foot --- The Ground Connection}\label{sec-ankle_foot}\label{sec:ankle_foot}

\subsection{Anatomy and Simplified Motion}\label{anatomy-and-simplified-motion}

The talocrural articulation principally contributes dorsiflexion--plantarflexion; subtalar and other foot articulations contribute coupled three-dimensional motion. Tarsometatarsal joints connect tarsals to metatarsals; metatarsophalangeal joints connect metatarsals to the proximal phalanges. They should not all be described as toe joints. A rigid-foot model removes deformation of this articulated structure.

Two effective ankle coordinates are a useful reduction, but treating their axes as exactly perpendicular sagittal and frontal hinges is an assumption. The foot can stay nearly stationary while ankle moments and contact forces change substantially. Conversely, heel lift, toe pivot and local slip can change support geometry even when a coarse video description calls the foot planted.

\subsection{Contact, Work and Deformation}\label{contact-work-and-deformation}

Ground contact needs a separate specification: contact locations or surfaces, normal compliance or an ideal nonpenetration constraint, friction, and rules for mode changes. A permanently welded foot, a flat sliding foot and unilateral frictional contact predict different reactions and admissible motion. Count only constraints active in the selected mode, and check whether their resulting reactions are physically feasible.

For the entire golfer--club system, ground reaction is an external wrench that changes momentum. A point force does zero power at an ideal stationary contact point. For a contact wrench, total power also includes its moment paired with angular velocity; an ideal fixed contact has zero translational and angular velocity. This does not make the ground irrelevant to power generation: support changes the motions and loads through which muscles do work. It does mean that force-plate force multiplied by center-of-mass velocity is not automatically the local work done at the foot--ground interface.

A deformable arch may store and return elastic energy. In an illustrative passive rotational spring--damper,

\[
\begin{gathered}U=\tfrac12 k\eta^2,\qquad \tau=-k\eta-c\dot\eta,\\ \tau\dot\eta=-\dot U-c\dot\eta^2,\end{gathered}
\]

for constant \(k>0\), \(c\geq0\). The spring returns previously stored energy; it does not create energy. If muscle activation changes effective stiffness, a variable parameter can exchange energy and must be included in the ledger. Identifying arch behavior requires deformation and load data, not simply a change in force-plate center of pressure.

\subsection{Choosing the Ankle/Foot Model}\label{choosing-the-anklefoot-model}

Choose a rigid foot for a task whose predictions remain adequate under that reduction. Add articulated toes, arch deformation or distributed contact when they materially change the measured output. ``Weight shift'' should be translated into observables such as individual foot forces, center of pressure, center-of-mass motion or changing contact mode; those quantities are not synonyms.

\section{Summary Table and Modeling Guidelines}\label{sec-summary_table}\label{sec:summary_table}

\begin{center}
\small
\begin{tabular}{@{}p{0.16\linewidth}p{0.34\linewidth}p{0.40\linewidth}@{}}
\toprule
Region & Candidate Reduction & What Still Needs Checking \\
\midrule
Hip & Spherical, three rotations & Frames, center, combined configuration limits \\
Knee & Hinge or coupled one-coordinate joint & Axial motion, translations, task sensitivity \\
Ankle/foot & Two effective rotations, rigid foot & Axis geometry, deformation and contact modes \\
Shoulder & Effective spherical arm root or explicit girdle & Moving center, GH versus thorax angles, muscle work \\
Elbow/forearm & Flexion and forearm rotation & Axis identification and dynamic coupling \\
Wrist & Two rotations with coupled force laws & Carpal reduction, stiffness and grip closure \\
\bottomrule
\end{tabular}
\end{center}

\subsection{Decision Flowchart}\label{decision-flowchart}

\begin{enumerate}
\def\labelenumi{\arabic{enumi}.}

\item
  Define the prediction: pose, club delivery, net moment, tissue load or clinical outcome. These require different evidence.
\item
  Specify bodies, frames, independent coordinates, inputs, constitutive laws and contact modes. A joint name alone is not a model.
\item
  Fit the relevant measurements and inspect residuals across the full task. Include measurement uncertainty and parameter identifiability.
\item
  Compare plausible model refinements using the same outputs and observations. Additional complexity is useful when it changes an important prediction and its parameters can be supported.
\item
  Validate on observations not used for fitting, and report which conclusions remain unchanged across model choices. Good kinematic fit does not independently validate muscle forces or injury predictions.
\end{enumerate}

\subsection{A Complete Counting Example}\label{a-complete-counting-example}

Consider an explicitly chosen tree: a floating pelvis (six freedoms), a three-rotation pelvis--thorax joint, two legs with hip/knee/ankle freedoms \(3+1+2\), and two arms whose shoulder, elbow and forearm, and wrist freedoms are \(3+2+2\). Attach the club rigidly to one hand. Before adding ground or second-hand closure, the local velocity dimension is

\[
\begin{aligned}d_{\rm tree}&=6+3+2(3+1+2)\\ &\quad+2(3+2+2)=35.\end{aligned}
\]

The rigid club attachment adds a body and a weld, with no new relative freedom. Let \(J_c\) stack the active ground and second-hand closure constraints on these tree coordinates. At a regular feasible configuration,

\[
d=35-\operatorname{rank}J_c.
\]

A six-component hand-pose closure removes six freedoms only if those constraints are independent on this model. Ground conditions and their rank must also be stated. The count is neither ``13 per arm'' nor a universal 30-DOF golfer. Independent muscle inputs, activation states and shaft deformation may change other dimensions without changing this particular rigid kinematic count.

\subsection{How the Choices Interact}\label{how-the-choices-interact}

Changing a joint's center or admissible motion changes hand and club kinematics, muscle moment arms and the feasible contact reactions. Changing passive stiffness changes forces and stored energy, without necessarily changing the number of coordinates. Changing the grip model changes loop closure and wrench transmission. Changing the chosen input from net torque to neural excitation adds activation and force-generation dynamics. These decisions interact, but none substitutes for the others.

For a task output \(y(q,\pi)\), with configuration \(q\) and geometry parameters \(\pi\), the first-order sensitivity is \(\delta y\simeq J_q\delta q+J_\pi\delta\pi\). The second term accounts for geometry changes even when the coordinate values stay fixed. Dynamic predictions require the associated changes to inertia, loads and constraints as well. At impact, event timing and the full club delivery state matter. A small anatomical angle can have a material effect on a sensitive task, while a large angle difference can be irrelevant to another output. State the output and calculate the sensitivity before calling a simplification negligible.

\section{Injury Biomechanics of the Golf Swing}\label{sec-ch22_injury}\label{sec:ch22_injury}

\subsection{From Motion to Tissue Load}\label{from-motion-to-tissue-load}

Anatomical modeling can help investigate loading mechanisms, but a kinematic limit is not a tissue-failure criterion. An ideal constraint is a mathematical restriction on motion; biological injury depends on the affected tissue, force or stress history, rate, exposure and individual capacity. Net joint moments, internal contact forces, muscle forces and local stresses are different quantities. Their uncertainties accumulate across the inference chain.

\subsection{Low Back Injury}\label{low-back-injury}

The lumbar region can experience compression, shear and torsional loading during a swing. Muscle forces and co-contraction can contribute substantially to internal loading; it is not determined just by ground force, visible torso rotation or the X-factor. The pelvis--spine--thorax chain is an articulated path, not by itself a closed kinematic loop.

Lim, Chow and Chae estimated L4--L5 loads using measured motion, force plates and an EMG-assisted optimization model in five male collegiate golfers. Their abstract reports mean peak compression above six body weights during the downswing, with peak anterior and medial shear in follow-through \citep{Lim2012LumbarLoads}. These are model estimates under that study's assumptions, not direct instrumented measurements inside professional golfers' spines or universal phase-specific limits.

An address angle or per-segment rotation limit cannot be used as a prospective golf injury threshold without evidence linking that defined measurement to subsequent outcomes. A model comparison should specify segment geometry, muscle recruitment, passive tissue laws, loading history and the actual outcome of interest. A calculated load change is a mechanical result; showing that a swing modification reduces pain or future injury requires separate clinical evidence.

\subsection{Wrist and Hand Injuries}\label{wrist-and-hand-injuries}

Club--ball contact force is not the force transmitted unchanged to a wrist. For a club considered as a separate system, with mass and center-of-mass velocity denoted below by the subscript C, the external-force balance in an inertial frame is

\[
F_{\rm hands}+F_{\rm ball}+F_{\rm other}=m_C\dot v_C.
\]

The angular-momentum balance is also needed, and a flexible shaft and distributed grip introduce further dynamics. Even when expressed about the same point in the same frame, two hand wrenches generally cannot be uniquely inferred from their sum. A contact impulse or peak ball force therefore does not establish the load on a particular wrist tendon, carpal structure or grip location.

Club angular velocity, shaft-axis spin and anatomical wrist joint rates are distinct measurements. A club rate cannot be relabeled as the rate of one anatomical wrist coordinate. To study a specific hand injury mechanism, measure the relevant contact history, estimate the appropriate tissue loads with uncertainty, and distinguish acute loading from repeated exposure. A list of possible injured structures does not validate a single causal mechanism for all of them.

\subsection{Medial Epicondylitis and Muscle Inference}\label{medial-epicondylitis-and-muscle-inference}

The common flexor--pronator origin near the medial epicondyle provides an anatomical reason to investigate repeated loads in that region. It does not establish a universal trail-arm cause or prove eccentric contraction from observed forearm rotation. Fiber shortening or lengthening depends on musculotendon geometry, activation and tendon behavior; surface EMG indicates electrical activity rather than force or contraction mode by itself.

Clinical labels and injury surveys should be used with their population, case definition and exposure denominator. A model of a symptomatic player can test mechanical alternatives, but it cannot diagnose tendon pathology or assign causation from one motion pattern.

\subsection{Control Authority Is a Separate Question}\label{control-authority-is-a-separate-question}

A high drift-to-control norm ratio is not an injury criterion and does not demonstrate that control is ineffective in the direction that matters. Finite-horizon task authority depends on input bounds, dynamics, initial state, contact constraints, available time and the selected output. An instantaneous zero-torque comparison changes a declared actuator input; it does not remove all activation, stored elastic energy or prior muscle work.

One can formulate a hypothesis that certain states and load histories leave fewer feasible low-load delivery options. Testing it requires a model with defined tissue-load quantities, an admissible control set, uncertainty and independent outcome evidence. Describing a swing as ``on autopilot'' supplies none of these. The mechanical investigation and the clinical conclusion should remain connected through explicit evidence rather than a metaphor.

\section{Chapter Exercises}\label{sec:ch22_exercises}

\begin{enumerate}
\def\labelenumi{\arabic{enumi}.}
\item
  \textbf{Knee Analysis.} Model an omitted 10-degree rotation about the tibial long axis. Does a tibial length of 40 cm determine ankle-point displacement? Calculate it for points on the axis and 5 cm from it. \textbf{Worked answer:} Length along the axis is irrelevant to this pure rotation. The two displacements are zero and \(2(0.05)\sin(5^\circ)\simeq0.00872\) m. Ankle orientation still changes in both cases. Translation, a different axis and chain closure require additional information.
\item
  \textbf{Hip Limits.} An individual's measured hip internal-rotation envelope is narrower than the model's default. Does that uniquely determine a compensatory spinal motion? \textbf{Worked answer:} No.~Express the measured restriction in compatible femur--pelvis frames and posture, then impose it on the coupled model. Foot orientation, pelvis motion, knee motion, task targets and other bounds determine whether feasible alternatives exist. A hip limit alone neither diagnoses FAI syndrome nor identifies a unique compensation.
\item
  \textbf{Shoulder Rhythm.} Under a deliberately planar 2:1 GH-to-scapular ratio, divide 70 degrees of total humerothoracic elevation into its components. Explain the limitation. \textbf{Worked answer:} GH contributes \(140/3\simeq46.67\) degrees and scapula \(70/3\simeq23.33\) degrees. This is arithmetic under the assumed same-axis ratio, not a prediction for a measured golf backswing. In three dimensions use relative rotation composition and measured or modeled girdle motion.
\item
  \textbf{Wrist Coupling.} State the dart-thrower's direction and calculate the energy in the illustrative stiffness matrix for \(\Delta q=(0.1,0.1)^T\) and \((0.1,-0.1)^T\) radians. \textbf{Worked answer:} The path goes from radial extension toward ulnar flexion. The energies are 0.01 J and 0.03 J. Direction-dependent resistance can be represented with two coordinates; this example supplies no empirical force or speed ranking.
\item
  \textbf{Elbow Carrying Angle.} Does a static 10-degree valgus alignment identify the full flexion axis? \textbf{Worked answer:} No.~It constrains an alignment measurement in the chosen posture, not all axis parameters. Define segment frames, obtain several flexion poses, fit an axis with residuals and uncertainty, and control for forearm rotation and marker motion. A muscle moment-arm estimate needs further attachment/path information.
\item
  \textbf{Grübler for a Golf Body.} Fix the pelvis to ground. Add distinct lumbar and thoracic bodies through two spherical joints, then two arms each with a spherical shoulder, hinge elbow and universal wrist. Weld a rigid club to one hand. Count before and after a generic full-rank second-hand pose closure. \textbf{Worked answer:} Merging pelvis with ground gives ten bodies including ground and nine tree joints, with freedoms \(3+3+2(3+1+2)=18\). The spatial count is \(6(10-1)-(54-18)=18\). A regular independent rank-six second-hand closure leaves twelve. This illustrative arm/spine model omits legs and forearm pronation; it is not the 35-coordinate model above. At singular or infeasible configurations the generic conclusion must be reconsidered.
\item
  \textbf{Model Refinement.} A model matches the backswing but releases a projected wrist angle too early. Does this identify insufficient wrist stiffness? \textbf{Worked answer:} No.~Alternative explanations include coordinate mismatch, initial velocity or activation error, grip or shoulder geometry, shaft dynamics, loading and timing. Compare predicted anatomical angles, hand paths and measured forces, then vary independently supported parameters and test held-out swings. A stiffness change that improves one fitted trace is not unique mechanistic identification.
\end{enumerate}
